\chapter{Logic for Compiler Reasoning}
\label{ch:logic-compiler-reasoning}

Replacing a division $x/x$ by the constant $1$ looks like elementary algebra.
Yet an optimizer must first ask whether $x$ can be zero, what division means in this IR, and whether evaluating $x$ has an observable effect.
A graph tells us where an instruction executes; logic tells us which claims justify changing it there.
The useful product of a proof is a precise connection between assumptions, execution, and the property we want to preserve.

We begin with pure expressions and mathematical integers, then make the boundary with machine operations explicit.
Statements in the small assertion language use deterministic, total expression evaluation unless stated otherwise.
This assumption deliberately excludes overflow, traps, memory faults, and side effects from expression evaluation.
Those features require additional premises, not a different interpretation of the word \enquote{proof}.
The logical foundations and program-proof rules developed here are treated systematically by Huth and Ryan~\cite{huth-2004-logic}.

\section{Facts at a program point}
\label{sec:ch40-facts}

Suppose a bounds check branches on $0\leq i$ and then on $i<n$.
At the array access, both comparisons have succeeded, provided neither operand has changed.
Write $p$ for $0\leq i$ and $q$ for $i<n$.
The access requires $p\land q$, not merely $p\rightarrow q$.
Table~\ref{tab:ch40-truth} shows why the distinction matters.

\begin{table}
\caption{An implication can hold even when the bounds-check conjunction fails.}
\label{tab:ch40-truth}
\centering
\begin{tabular}{ccccc}
\toprule
$p$ & $q$ & $\neg p$ & $p\land q$ & $p\rightarrow q$ \\
\midrule
F & F & T & F & T \\
F & T & T & F & T \\
T & F & F & F & F \\
T & T & F & T & T \\
\bottomrule
\end{tabular}
\end{table}

A \term{propositional formula} combines atomic propositions using connectives such as negation $\neg$, conjunction $\land$, disjunction $\lor$, and implication $\rightarrow$.
A \term{valuation} assigns each atom either true or false.
Evaluation extends this assignment through the formula's syntax tree.
Conjunction requires both operands to be true; disjunction requires at least one; negation reverses the truth value.
The formula $A\rightarrow B$ means $\neg A\lor B$.
It is false only when $A$ is true and $B$ is false.
The constants $\top$ and $\bot$ denote truth and falsity.

Logical implication describes truth under an assignment, not temporal execution or causation.
Similarly, $A\land B$ does not specify which operand a program evaluates first.
If a source expression uses short-circuit evaluation, that control flow must be represented separately before logical formulas can describe its results.
In particular, a logical formula containing a division is not made operationally safe merely by writing a nonzero test to its left.

A formula is \term{satisfiable} if some valuation makes it true, \term{valid} if every valuation does, and \term{unsatisfiable} if none does.
The formula $p\lor\neg p$ is valid; $p\land\neg p$ is unsatisfiable; $p\land q$ is satisfiable but not valid.
These are properties of formulas, not three possible truth values of an individual valuation.

Treating comparisons as independent atoms can lose arithmetic information.
For example, the propositional assignment that makes both $x<0$ and $x\geq0$ true has no corresponding integer state.
A propositional proof valid under all assignments remains useful under this abstraction, but a propositional counterexample might be arithmetically impossible.
Predicate semantics will recover the relationship between the atoms.

Figure~\ref{fig:ch40-guard-facts} attaches assertions to the successful edges of the bounds check.
The fact belongs to a particular version of $i$ and $n$.
If an intervening assignment changes $i$, the optimizer cannot reuse the old comparison without proving a relationship between the old and new values.
Dominance identifies a must-pass-through block; it does not by itself prove that a value tested there is still unchanged.

\begin{figure}
\centering
\begin{tikzpicture}[node distance=12mm and 23mm]
\node[basic block] (lower) {test $0\leq i$};
\node[basic block, below=of lower] (upper) {test $i<n$};
\node[basic block, below=of upper] (access) {access $A[i]$};
\node[basic block, right=of upper] (fail) {reject};
\draw[cfg edge] (lower) -- node[left, font=\small] {$p$} (upper);
\draw[cfg edge] (upper) -- node[left, font=\small] {$p\land q$} (access);
\draw[cfg edge] (lower.east) -| node[above, pos=0.35, font=\small] {$\neg p$} (fail.north);
\draw[cfg edge] (upper) -- node[above, font=\small] {$p\land\neg q$} (fail);
\end{tikzpicture}

\caption{Successful test edges accumulate facts about unchanged operands before an access.}
\label{fig:ch40-guard-facts}
\end{figure}

\section{Meaning and derivation are different objects}
\label{sec:ch40-entailment}

Let $\Gamma$ be a finite collection of assumptions.
We write $\Gamma\models A$ for \term{semantic entailment}: every interpretation satisfying all members of $\Gamma$ also satisfies $A$.
For propositional formulas, interpretations are valuations.
An assignment that makes every assumption true and $A$ false is a counterexample to the entailment.
Thus $\Gamma\models A$ exactly when
\[
  \left(\bigwedge_{G\in\Gamma}G\right)\land\neg A
\]
is unsatisfiable; the empty conjunction is $\top$.
This reformulation is the bridge to solver-based checking in Chapter~\ref{ch:automated-reasoning-limits}.

The notation $\Gamma\vdash A$ means \term{derivability}: there exists a finite proof of $A$ from $\Gamma$ in a specified proof system.
Entailment quantifies over meanings; derivability quantifies over proof objects.
For example, the following rules introduce a conjunction, extract one of its components, and apply an implication:
\[
  \frac{\Gamma\vdash A\qquad\Gamma\vdash B}{\Gamma\vdash A\land B}
  \qquad
  \frac{\Gamma\vdash A\land B}{\Gamma\vdash A}
  \qquad
  \frac{\Gamma\vdash A\qquad\Gamma\vdash A\rightarrow B}
       {\Gamma\vdash B}.
\]
Each rule is a finite instruction for constructing a larger proof from smaller proofs.
An assumption rule permits $\Gamma\vdash A$ when $A\in\Gamma$.

Take $\Gamma=\{p,\ p\rightarrow q,\ q\rightarrow r\}$.
The assumption rule provides $p$ and $p\rightarrow q$; implication elimination yields $q$.
A second application, using $q\rightarrow r$, yields $r$.
Conjunction introduction then yields $p\land r$.
This is a derivation with named dependencies, rather than a claim that the conclusion merely looks plausible.

A proof system is sound when $\Gamma\vdash A$ implies $\Gamma\models A$.
It is complete for the chosen semantics when the converse holds.
The few displayed rules are not a complete system; standard classical propositional calculi add rules for the other connectives and are sound and complete~\cite{huth-2004-logic}.
Completeness does not promise a short proof or a fast search procedure.
Nor does soundness of a calculus establish that a compiler's encoding of an instruction is correct.
The statement proved and its connection to the implementation are separate obligations.

Inconsistent assumptions entail every conclusion: no interpretation satisfies them, so none can refute the entailment.
An optimization checker that assumes both $x=3$ and $x=4$ can therefore \enquote{prove} a false rewrite conditionally on those assumptions.
This may correctly describe an unreachable path, or it may expose an encoding bug.
Checking whether assumptions are satisfiable is a useful diagnostic; it is not a replacement for checking that they describe reachable program states soundly.

\section{Predicates, domains, and binding}
\label{sec:ch40-quantifiers}

The statement $i<n$ has structure that an atom named $q$ hides.
In \term{predicate logic}, terms denote domain elements, predicates denote relations on those elements, and quantifiers bind variables.
An interpretation $\mathcal M$ supplies a nonempty domain, meanings for constants, meanings for function symbols, and meanings for predicate symbols.
An environment $\rho$ supplies values for free variables.
We write $\mathcal M,\rho\models A$ when formula $A$ is true in that interpretation and environment.
For arithmetic assertions here, the domain is $\mathbb Z$, addition has its usual meaning, and $<$ is the usual integer ordering.

The universal assertion $\forall x.\,A$ holds when $A$ holds for every replacement of $\rho(x)$ by a domain element.
The existential assertion $\exists x.\,A$ holds when at least one replacement makes $A$ true.
A variable occurrence is \term{bound} when it is within the scope of a quantifier for that variable; otherwise it is \term{free}.
A formula with no free variables is a sentence.
An assertion about a running program normally has free variables: its truth depends on the current state.

Parentheses specify scope.
The formula $(\forall x.\,P(x))\lor Q(x)$ has a free occurrence of $x$ in $Q(x)$.
The formula $\forall x.\,(P(x)\lor Q(x))$ binds both occurrences.
Renaming a bound variable consistently, without capturing other variables, changes no meaning.
Using distinct internal identifiers for binders is therefore safer than identifying variables only by their printed names.

Quantifier order records dependency.
Consider two uses $u_1,u_2$ and two definitions $d_1,d_2$, where the relation $R(u,d)$ consists exactly of $(u_1,d_1)$ and $(u_2,d_2)$.
Then
\[
 \forall u\in\{u_1,u_2\}.\,\exists d\in\{d_1,d_2\}.\,R(u,d)
\]
is true: each use has a suitable definition.
But the formula with $\exists d$ before $\forall u$ is false: no single definition serves both uses.
The bounded notation $\forall x\in S.\,A$ abbreviates $\forall x.\,(x\in S\rightarrow A)$; existential restriction uses conjunction instead.
Reordering quantifiers can silently change a per-use obligation into a common-witness obligation.

To instantiate a universally quantified theorem, substitute a term for its bound variable while respecting other binders.
To prove a universal claim, reason about an arbitrary fresh variable, with no assumption that restricts it to a convenient value.
To use an existential assumption, introduce a fresh name for its witness within a local argument; a conclusion that escapes that argument must not depend on which fresh name was chosen.
These freshness conditions prevent a proof from treating an unknown value as a particular constant.

\subsection{Substitution must avoid capture}
\label{sec:ch40-substitution}

Write $Q[e/x]$ for replacement of the free occurrences of $x$ in $Q$ by expression $e$.
Bound occurrences of $x$ are untouched.
Any binder that would capture a free variable of $e$ must first be renamed to a fresh variable.
This operation is \term{capture-avoiding substitution}.

For example, let $Q=\forall y.\,x\leq y$ and substitute the free variable $y$ for $x$.
Naive text replacement produces $\forall y.\,y\leq y$, which is always true.
Correct substitution first renames the bound variable and produces $\forall z.\,y\leq z$.
Over the domain $\{0,1\}$ with the usual ordering and free $y=1$, the correct result is false because $z=0$ refutes it.
The naive result remains true.
The difference is a change in meaning, not typography.

Compiler-generated assertions often avoid source-level name clashes by giving every program value and every logical binder a unique identity.
Printing can recover readable names afterwards.
Nevertheless, a substitution routine must work on binding structure, not substrings: replacing the letter \code{x} in an identifier is not logical substitution.

\section{Proof by induction follows the object being proved}
\label{sec:ch40-induction}

A compiler correctness argument often covers infinitely many expressions or executions using a finite collection of cases.
The induction measure explains why those cases suffice.
For a syntax-directed transformation, use \term{structural induction}: prove the property for leaves, then prove it for each constructor assuming it holds for the constructor's children.
For a proof checker, induct on the derivation tree.
For a loop invariant, induct on the number of completed iterations.
These are different objects, so their induction hypotheses are different.

As a small structural example, consider pure integer expressions built from constants, variables, addition, and negation.
A recursive simplifier processes children first and removes a double negation when its simplified child is itself a negation.
For a variable or constant, evaluation is unchanged.
For addition, the two induction hypotheses say that the simplified operands have the same values as the original operands; adding equal values gives equal results.
For negation, either the constructor remains and the induction hypothesis suffices, or a simplified operand $-t$ is exposed and the result is $t$.
In the latter case the original negated expression evaluates to $-(-\llbracket t\rrbracket_\rho)=\llbracket t\rrbracket_\rho$.
This last equality uses the specified arithmetic semantics.
It would not justify silently ignoring an exception from an intermediate operation in a trapping arithmetic language.

Soundness of an inference system follows the same pattern.
Assumption leaves are true in every interpretation satisfying $\Gamma$.
At an implication-elimination node, the induction hypotheses establish both $A$ and $A\rightarrow B$; their truth conditions establish $B$.
Check each remaining rule similarly.
Induction then covers every finite derivation, not only the derivation used as an example.

Induction is also what justifies the assignment rule below.
For a total pure expression $e$, let $v=\llbracket e\rrbracket_\rho$ and let $\rho[x\mapsto v]$ update just $x$.
The substitution property is
\[
  \mathcal M,\rho\models Q[e/x]
  \quad\Longleftrightarrow\quad
  \mathcal M,\rho[x\mapsto v]\models Q.
\]
Its proof proceeds through the terms and formula constructors, with binder renaming supplying the quantifier case.
An implementation of assertion rewriting relies on this property every time it crosses an assignment backwards.

\section{Assertions before and after execution}
\label{sec:ch40-hoare}

A \term{Hoare triple} $\{P\}\ C\ \{Q\}$ connects a precondition $P$, a command $C$, and a postcondition $Q$.
For \term{partial correctness}, it says that every terminating execution of $C$ starting in a state satisfying $P$ ends in a state satisfying $Q$.
For the total-expression command language used here, there are no stuck executions.
When extending the language with faults, state explicitly whether the specification also requires their absence; a fault must not disappear merely because it fails to produce an ordinary final state.

\term{Total correctness} additionally requires termination from every state satisfying $P$.
Partial correctness alone validates $\{\top\}\ \text{while true do skip}\ \{\bot\}$: there is no terminating execution that could violate the postcondition.
It does not validate the corresponding total-correctness claim.
This distinction matters when considering transformations that delete loops or move effects across potentially nonterminating computations.

For assignment, the substitution property gives
\[
  \{Q[e/x]\}\ x:=e\ \{Q\}.
\]
If the desired postcondition after $x:=x+2$ is $x\geq7$, the required assertion before it is $x+2\geq7$, or $x\geq5$.
The occurrences of $x$ in $x+2$ denote the old value.
Substitution is performed once; the expression inserted for $x$ is not recursively rewritten by the same substitution.

Two additional rules connect local proofs:
\[
 \frac{\{P\}\ C_1\ \{R\}\qquad\{R\}\ C_2\ \{Q\}}
      {\{P\}\ C_1;C_2\ \{Q\}},
 \qquad
 \frac{P\models P'\quad\{P'\}\ C\ \{Q'\}\quad Q'\models Q}
      {\{P\}\ C\ \{Q\}}.
\]
The second is the consequence rule: strengthen the precondition or weaken the postcondition.
Reversing either implication is unsound.
A proof requiring $x\geq5$ can accept callers promising $x\geq8$, not callers promising only $x\geq0$.

\subsection{Computing obligations backwards}
\label{sec:ch40-backward}

For loop-free commands with total expressions, the \term{weakest precondition} $\operatorname{wp}(C,Q)$ is exactly the set of initial states whose execution finishes in $Q$.
It is weakest in the implication order: every sufficient precondition entails it.
Algorithm~\ref{alg:ch40-wp} computes a formula for that set by recursion over structured commands.
The command \code{skip} changes no state, and a conditional evaluates its guard before executing exactly one arm.

\begin{algorithm}
\caption{Construct a weakest precondition for a loop-free command.}
\label{alg:ch40-wp}
\begin{pseudocode}
\Require Command $C$ built from skip, assignment, sequence, and conditionals
\Require Postcondition $Q$; pure, total expressions and guards
\Ensure Formula true exactly in the states from which $C$ establishes $Q$
\Function{WP}{$C,Q$}
  \If{$C$ is skip}
    \State \Return $Q$
  \ElsIf{$C$ is $x:=e$}
    \State Rename binders in $Q$ that would capture free variables of $e$
    \State \Return $Q[e/x]$ using structural substitution
  \ElsIf{$C$ is $C_1;C_2$}
    \State $R\gets\Call{WP}{C_2,Q}$
    \State \Return $\Call{WP}{C_1,R}$
  \Else \Comment{$C$ is if $b$ then $C_t$ else $C_f$}
    \State $R_t\gets\Call{WP}{C_t,Q}$
    \State $R_f\gets\Call{WP}{C_f,Q}$
    \State \Return $(b\rightarrow R_t)\land(\neg b\rightarrow R_f)$
  \EndIf
\EndFunction
\end{pseudocode}
\end{algorithm}

For $a:=a+2;\ b:=3a$ with goal $b\geq12$, crossing the second assignment gives $3a\geq12$.
Crossing the first gives $3(a+2)\geq12$, equivalent over integers to $a\geq2$.
Figure~\ref{fig:ch40-backward} shows the forward state changes beside these backward obligations.
At $a=2$, the final $b$ is $12$; at $a=1$, it is $9$, so weakening the precondition to $a\geq1$ would be wrong.

\begin{figure}
\centering
\begin{tikzpicture}[node distance=14mm]
\node[basic block] (first) {$a:=a+2$};
\node[basic block, below=of first] (second) {$b:=3a$};
\node[basic block, below=of second] (end) {finish};
\draw[cfg edge] (first) -- (second);
\draw[cfg edge] (second) -- (end);
\node[right=20mm of first, align=center] (p) {$3(a+2)\geq12$};
\node[right=20mm of second, align=center] (r) {$3a\geq12$};
\node[right=20mm of end, align=center] (q) {$b\geq12$};
\draw[cfg edge] (q) -- (r);
\draw[cfg edge] (r) -- (p);
\node[above=5mm of first, font=\small] {execution};
\node[above=5mm of p, font=\small] {obligation};
\end{tikzpicture}

\caption{Execution moves down the command sequence while substitution moves the obligation up it.}
\label{fig:ch40-backward}
\end{figure}

For a conditional example, execute $x:=-x$ when $x<0$, otherwise execute $x:=x+2$, and then execute $y:=3x$.
The goal $y\geq6$ first becomes $x\geq2$.
The two arms require $-x\geq2$ and $x+2\geq2$, respectively, giving
\[
  (x<0\rightarrow x\leq-2)\land
  (x\geq0\rightarrow x\geq0).
\]
Over integers this is equivalent to $x\neq-1$.
The input $-2$ produces $y=6$, the input $-1$ produces $y=3$, and the input $0$ produces $y=6$.
The guard is evaluated in the incoming state, so it must not be replaced by the post-assignment value of $x$.

Let $n$ be the number of command nodes and $L$ the total number of syntax nodes inspected or allocated by all substitution and formula-construction steps.
With linear-time capture avoidance and constant-time node operations, the running time is $O(n+L)$; recursive control storage is $O(n)$ in the worst case, in addition to formula storage.
This is an output-sensitive statement, not a claim that $L$ is polynomial in $n$.
Repeated substitution through $x:=x+x$ can double a tree representation at each step.
Shared expression graphs reduce duplication, but substitution caches must include the substitution environment, not only the original node.
The algorithm constructs a formula; it does not decide arithmetic validity or simplify that formula automatically.

\subsection{Loops need an inductive assertion}
\label{sec:ch40-loop-proof}

Backward recursion cannot unfold a loop indefinitely.
Instead choose a \term{loop invariant}, an assertion $I$ that holds at the loop head initially and after every completed iteration.
For a loop with guard $b$, body $C$, and desired exit property $Q$, prove initialization, preservation, and exit:
\[
  P\models I,\qquad
  \{I\land b\}\ C\ \{I\},\qquad
  I\land\neg b\models Q.
\]
Here $P$ describes the state just before the first guard evaluation.
The preservation triple establishes the induction step for every state satisfying $I\land b$, not only for states reached in one sample execution.

Take an unchanged integer $n\geq0$, initialize $i:=0$ and $t:=0$, then repeat $t:=t+3;\ i:=i+1$ while $i<n$.
Choose $I\equiv 0\leq i\leq n\land t=3i$.
Initialization satisfies $I$.
Under $I\land i<n$, integer discreteness gives $i+1\leq n$, and the next values satisfy $t'=t+3=3(i+1)=3i'$.
At exit, $i\leq n$ and $i\geq n$ give $i=n$, hence $t=3n$.
Table~\ref{tab:ch40-loop-trace} checks the arithmetic for $n=4$.

\begin{table}
\caption{The loop invariant and decreasing variant at each guard evaluation for $n=4$.}
\label{tab:ch40-loop-trace}
\centering
\begin{tabular}{rrrrc}
\toprule
$i$ & $t$ & $3i$ & $n-i$ & $i<n$ \\
\midrule
0 & 0 & 0 & 4 & T \\
1 & 3 & 3 & 3 & T \\
2 & 6 & 6 & 2 & T \\
3 & 9 & 9 & 1 & T \\
4 & 12 & 12 & 0 & F \\
\bottomrule
\end{tabular}
\end{table}

For total correctness, use the \term{loop variant} $n-i$.
It is a nonnegative integer at the head and decreases by one on each taken iteration; the body itself terminates.
There is no infinite strictly descending sequence of nonnegative integers, so the loop terminates.
A real-valued quantity that merely decreases is insufficient: $1,1/2,1/4,\ldots$ decreases forever while remaining positive.
Termination requires a well-founded descent, not just a less-than comparison between successive values.

\section{Turning a proof into an optimization precondition}
\label{sec:ch40-optimization}

For pure total expressions $e$ and $e'$, a local rewrite under fact $P$ requires
\[
  \forall\rho.\ P(\rho)\rightarrow
  \llbracket e\rrbracket_\rho=\llbracket e'\rrbracket_\rho.
\]
The quantification is over all admissible input states, not just inputs seen in testing.
Substitution of equal expressions into a larger pure expression preserves its value because each enclosing operator is a function of its operands.
This compositional argument stops being sufficient when evaluation order, effects, or exceptional results become observable.

For example, replacing $x/x$ by $1$ is justified on nonzero mathematical integers with exact integer division of equal operands.
It is also an equality of the corresponding defined integer operations for nonzero machine integers when that operation has no other exceptional case on equal operands.
It is not an unconditional floating-point identity: zero, infinities, and NaNs need separate treatment.
The useful compiler rule includes its type and definedness premises, rather than relying on the spelling of the operator.

Code motion needs an execution premise as well as a value equality.
Suppose a division executes only on the true arm of a branch.
Moving it before the branch introduces an evaluation on the false arm.
Even if neither operand changes, the new evaluation can fault when the denominator is zero there.
A proof must establish safety on every newly exposed path, or show that the IR semantics permits the speculation.
Likewise, an address equality does not justify merging loads if memory may change between them.

Facts supplied by analysis are premises in these proofs.
An interval result $0\leq x\leq14$ can justify a comparison rewrite for four-bit unsigned addition, while $0\leq x\leq15$ cannot justify the same claim.
An implementation should record the fact's program point, operand identities, arithmetic semantics, and invalidation conditions.
This gives a concrete debugging order: inspect the premises, inspect the semantic rule, then inspect the transformation that applied it.

\section{Exercises}
\label{sec:ch40-exercises}

Solutions and finite-example checks are in \code{example-code/logic-checks/}.
The checks supplement the general arguments; bounded enumeration is not a proof over all integers.

\begin{enumerate}
\item\label{ex:ch40-entailment}
Let $\Gamma=\{p\lor q,\ p\rightarrow r\}$.
Give a valuation refuting $\Gamma\models r$, then show that adding $\neg q$ makes the entailment valid.

\item\label{ex:ch40-capture}
Compute the capture-avoiding substitution $Q[y/x]$ for $Q=\exists y.\,x<y$.
Over $\mathbb Z$ with free $y=5$, compare its truth with that of the naive text replacement.

\item\label{ex:ch40-sequence}
Over integers, compute the weakest precondition of $a:=a+2;\ b:=3a$ for $b\geq15$.
Give one boundary input that satisfies it and one that does not.

\item\label{ex:ch40-branch}
When $x<0$, execute $x:=-x$; otherwise execute $x:=x+2$.
Then execute $y:=3x$.
Compute the weakest precondition for $y\geq9$ over integers, retaining the guard constraints until the final simplification.

\item\label{ex:ch40-invariant}
Change the loop in Section~\ref{sec:ch40-loop-proof} to increment $t$ by $5$ instead of $3$.
Give an invariant and a variant proving total correctness for unchanged $n\geq0$, and compute the exit state for $n=3$.

\item\label{ex:ch40-speculation}
An IR division traps on zero.
A program evaluates $12/d$ only when a Boolean input $b$ is true.
An optimizer moves the division before the branch, assuming only $b\rightarrow d\neq0$.
Give a counterexample and a sufficient stronger premise for this motion, assuming the division is otherwise pure and total.
\end{enumerate}

\section{Further reading}
\label{sec:ch40-reading}

Huth and Ryan~\cite{huth-2004-logic} connect propositional and predicate proof systems to program assertions; their discussions of substitution and partial versus total correctness explain the side conditions that a proof-producing compiler must preserve.
Boolos, Burgess, and Jeffrey~\cite{boolos-2007-computability} develop the relationship between formal derivations and semantic consequence, providing the background for understanding why a complete calculus need not yield a terminating decision procedure.
